\subsection{Diseño de Excepciones Propias}

\subsubsection{Regla: constructores de una excepción propia}

\textbf{Regla:} Una excepción propia hereda de \texttt{Exception}, se declara \texttt{sealed}, termina en \texttt{Exception} (Sección 2.14.1), fija su mensaje en inglés dentro del constructor y expone sus datos mediante propiedades de solo lectura. Si no lleva datos, declara los tres constructores estándar que exige la regla de análisis \texttt{CA1032} (Microsoft, 2016b). Si lleva datos obligatorios, declara solo los constructores que los reciben.

\textbf{Descripción:} Un constructor sin parámetros dejaría una excepción con datos sin asignar, lo que contradice la Sección 2.6.3. Como la clase es \texttt{sealed}, ninguna clase derivada necesita los constructores restantes.

\textbf{Código correcto:}
\begin{lstlisting}[style=csharpstyle]
public sealed class TileOccupiedException : Exception
{
    public TileOccupiedException(Tile tile)
        : base($"The tile is already occupied. Row={tile.Row} Column={tile.Column}")
    {
        Tile = tile;
    }

    public Tile Tile { get; }
}
\end{lstlisting}

\textbf{Código incorrecto:}
\begin{lstlisting}[style=csharpstyle]
public class TileOccupiedException : Exception
{
    public TileOccupiedException()
    {
    }

    public Tile? Tile { get; set; }
}
\end{lstlisting}
